%Datum und Uhrzeit
\usepackage{datetime}

%Encoding
\usepackage[utf8]{inputenc}
\usepackage{lmodern}

%Graphikpakete
\usepackage{graphicx}
\usepackage[usenames,dvipsnames]{xcolor}
\usepackage{tikz}
\usetikzlibrary{arrows, snakes, backgrounds}
\usepackage{wrapfig}
\usepackage{dirtree}

%------------------------------------------------------------------
%% THI - Farbdefinitionen
%------------------------------------------------------------------

% Farben der THI
\definecolor{THIblue}{rgb}{0.0078,0.1176,0.4705}
\definecolor{AllianzBlue}{rgb}{0.07,0.20,0.53}
%%-------------Allgemeine Definitionen----------------------------------
% Farbige Aufwertung der berschriften
\addtokomafont{chapter}{\color{AllianzBlue}}
\addtokomafont{section}{\color{AllianzBlue}}
\addtokomafont{caption}{\color{AllianzBlue}}
\addtokomafont{subsection}{\color{AllianzBlue}}
\addtokomafont{subsubsection}{\color{AllianzBlue}}
\setkomafont{captionlabel}{\color{AllianzBlue}}


%Zitierumgebung
%\usepackage{cite}% Zitieren
\usepackage[ngerman]{babel}
\usepackage[backend=biber,]{biblatex}
%\usepackage{bibgerm}% Literatur in Deutscher DIN
\usepackage[babel,german=quotes]{csquotes}
\bibliography{ref/ref_liste} % Pfad und Datei der Ref-Datenbank

%URL-Umgebung
\usepackage{url}
\renewcommand{\UrlFont}{\small\color{THIblue}}

%Mathematik
\usepackage{amsmath}
\usepackage{amssymb}
\usepackage{mathtools}

%Tabellen
\usepackage{longtable}
\usepackage{multirow}

%Quellcode
\usepackage{listings}
\lstset{literate=%
    {Ö}{{\"O}}1
    {Ä}{{\"A}}1
    {Ü}{{\"U}}1
    {ß}{{\ss}}1
    {ü}{{\"u}}1
    {ä}{{\"a}}1
    {ö}{{\"o}}1
    {~}{{\textasciitilde}}1
}

\lstset{% general command to set parameter(s)
	basicstyle=\small\ttfamily,%\small, % print whole listing small
	keywordstyle=\color{THIblue}\bfseries,%\underbar,% underlined boldblack keywords
	identifierstyle=, % nothing happens
	commentstyle=\color{Aquamarine}, % white comments
	stringstyle=\color{Cerulean}\ttfamily, % typewriter type for strings
	showstringspaces=false, % no special string spaces
	frame=single,
	framexleftmargin=1.5em,
	xleftmargin=2em,
	numbers=left,
	numberstyle=\color{gray},
	numbersep=5pt,
	captionpos=b,
	breaklines=true}

\lstdefinelanguage{conf}
{
  % list of keywords
  basicstyle=\normalsize\ttfamily, %lineend
  sensitive=false, % keywords are not case-sensitive
  morecomment=[l]{\#}, % l is for line comment
  morecomment=[s]{/*}{*/}, % s is for start and end delimiter
  morestring=[b]" % defines that strings are enclosed in
}

\lstdefinelanguage{tikz}
{
  % list of keywords
  morekeywords={
	above,
	align,
  	begin,
	below,
  	caption,
	draw,
	end,
	every,
	fill,
	label,
  	line,
	minumum,
	node,
	rectangle,
	right,
	scale,
	size,
  	style,
	of,
  	width,
  	xshift,
  	yshift
  },
  sensitive=false, % keywords are not case-sensitive
  %alsodigit={0482},
  keywordstyle=\color{THIblue}\bfseries,
  morecomment=[l]{//}, % l is for line comment
  morecomment=[s]{/*}{*/}, % s is for start and end delimiter
  morestring=[b]" % defines that strings are enclosed in double quotes
}

%Microtypographie
\usepackage{microtype}

%Kopf und Fußzeilen
\usepackage{scrpage2} 	% Kopf & Fußzeile im KOMA Stil
\pagestyle{scrheadings}	% Aktiviert Verwendung vordefinierter Kolumnentitel
\clearscrheadfoot 		% alle Standard-Werte und Formatierungen löschen
\setkomafont{pagehead}{\scshape}	% Schriftart in Kopfzeile, \scshape = Kapitelchen
\automark[chapter]{section} % [linke Seite]{rechte Seite}
%\ohead{\def\pagestyle{PDTS}{\hrulefill\includegraphics[width = 6cm]{bilder/thi_logo_quer_cropped}}}
\ohead{\includegraphics[width = 6cm]{bilder/thi_logo_quer_cropped}}

%\ihead{\textsc{Abschlussarbeit}}
\ihead{\headmark}

%\setheadwidth[0pt]{textwithmarginpar}
\ofoot{\vspace{-0.3cm} \pagemark} 						
\ifoot{\vspace{-0.3cm} Dominik Gunther Florian Schlecht} 
				
%\setheadtopline{.4pt}				
\setheadsepline{.2pt}
\setfootsepline{.4pt}	% Trennlinie Fußzeile und Textkörper

%------------------------------------------------------------------
%% Längenanpassungen
%------------------------------------------------------------------
\setlength{\headsep}{10mm}		% Textabstand zur Kopfzeile
\setlength{\footskip}{15mm}		% Abstand zur Fußzeile
\setlength{\parindent}{0em}		% Einzug nach Absatz

%%-------------Allgemeine Definitionen----------------------------------
% Farbige Aufwertung der berschriften
\addtokomafont{chapter}{\color{THIblue}}
\addtokomafont{section}{\color{THIblue}}
\addtokomafont{caption}{\color{THIblue}}
\addtokomafont{subsection}{\color{THIblue}}
\addtokomafont{subsubsection}{\color{THIblue}}
\setkomafont{captionlabel}{\color{THIblue}}

%Hyperref
\usepackage[
		pdftex,
		linkcolor=THIblue,			% Farbe der Verlinkung
		linktocpage=true,			% Im TOC wird Seitenzahl verlinkt(true),bzw. Text(false)
		colorlinks=true,			% 'true' keinen Kasten um Link
		citecolor=THIblue,
%		pdfhighlight=/P,
%		bookmarks,
%		hyperfigures=true,
%		citebordercolor={0 0 1},
%		linkbordercolor={0 0 1},
%		menubordercolor={0 0 1},
%		backref=true,
%		pagebackref=true,
%		bookmarksopen,
%		bookmarksnumbered,
%		pdfpagelabels=false,
%		pdfstartpage=1,
%		pdfstartview=Fit,			% Modus beim Öffnen (Fit = An Seitengröße anpassen)
		pdftitle={Umgehen von USB-Deskriptor basierten USB-Policies am Beispiel einer virtuellen Umgebung},
		pdfauthor={Dominik Schlecht},
%		pdfstartview=Fit,
%		pdfdisplaydoctitle=true,
%		plainpages=false
			]{hyperref}  

%------------------------------------------------------------------
% Wichtige Definition für Aufteilung von Formelverzeichnis und Abkürzungsverzeichnis
%------------------------------------------------------------------

% Nomenklaturverzeichnis, Formelzeichenliste Anpassungen für nomenclature: damit lassen sich zwei getrennte Symbolverzeichnisse anlegen, ziehmlich cool!
\usepackage[intoc,compatible,german]{nomencl}	
		\renewcommand{\nomgroup}[1]{	\ifthenelse{\equal{#1}{A}}{\item[{\normalfont\sffamily\bfseries\LARGE\textcolor[rgb]{0,0.112,0.47}{Abkürzungen{\phantom{\Huge $\frac{\frac{\frac{A}{a}}{a}}{\frac{a}{a}}$}}}}]}{	\ifthenelse{\equal{#1}{A}}{\item[{\normalfont\sffamily\bfseries\LARGE\textcolor[rgb]{0,0.112,0.47}{Formelzeichen{\phantom{\Huge $\frac{A}{\frac{a}{a}}$}}}}]}{}}}
		
%------------------------------------------------------------------
%% Anpassung von Abständen, Längen vom Nomenclaturverzeichnis (Abkürzungs- und Formelverzeichnis)
%------------------------------------------------------------------

% Abstand zwischen Einträgen im Symbolverzeichnis (-\parsep = 0)
\setlength{\nomitemsep}{-\parsep} 
\setlength{\nomlabelwidth}{5em}	
\renewcommand{\nomlabel}[1]{#1 \dotfill}  % Punkte im zwischen Nummer und Kapiteleintrag

% Sub-Figures mit Caption
\usepackage{subcaption}

% Paragraph as Subsubsubsection
\usepackage{titlesec}
\titleformat{\paragraph}
{\normalfont\normalsize\bfseries}{\theparagraph}{1em}{}
\titlespacing*{\paragraph}
{0pt}{3.25ex plus 1ex minus .2ex}{1.5ex plus .2ex}
